% TOP_PROBLEM: {"id":30007738,"problem_number":"LOCAL-30007738","title":"Policies and completed fertility","category_id":21,"category":{"id":21,"name":"economics","display_name":"Economics","description":"Open economic theory statements, published research agendas, and proposed scoped specifications; question provenance and historical priority are qualified per record.","slug":"economics","order_index":21,"created_at":"2026-10-08T00:00:00Z"},"status":"provenance_pending","external_url":null,"legacy_ids":[],"published":false,"statement":"Determine whether a defined childcare subsidy increases births completed by age forty-five rather than merely changing birth timing, and measure its resource cost.","economics":{"original_id":227,"field":"Education, families and demography","kind":"empirical","formalization_basis":"proposed_specification","source_status":"not_verified","priority_status":"not_established","priority_note":"No readable primary passage explicitly stating this precise open question was verified. A supporting paper, abstract, or topic citation is insufficient to establish origin or unresolved status.","earliest_verified_reference":null,"current_reference":{"authors":["Melissa S. Kearney","Phillip B. Levine"],"title":"Why Is Fertility So Low in High-Income Countries?","year":2026,"url":"https://www.aeaweb.org/articles?id=10.1257/jel.20261786","doi":"10.1257/jel.20261786","locator":"Abstract","evidence_summary":"The published review discusses explanations for falling fertility; the proposed counterfactual decomposition and completed-fertility policy contrast are our specifications."},"formalization_note":"Proposed scoped research specification. The original catalogue prompt was a synthesis, and no canonical conjecture or verified original formulation is claimed. The supporting source, when retained, establishes context or existing results rather than this exact open question."},"statusQualification":"Provenance pending: an explicit published statement of this exact open question was not verified. This is a catalogue-authored candidate specification, not a certified open conjecture. Proposed scoped research specification. The original catalogue prompt was a synthesis, and no canonical conjecture or verified original formulation is claimed. The supporting source, when retained, establishes context or existing results rather than this exact open question.","statusReviewedAt":"2026-10-08"}
% FIRST_POSTED: 2026-10-08
% ENTRY_KIND: statement_only
% !TeX program = lualatex
\documentclass[11pt]{article}
\usepackage{fontspec}
\usepackage[a4paper,margin=25mm]{geometry}
\usepackage{amsmath,amssymb,mathtools}
\usepackage{xurl}
\usepackage[hidelinks]{hyperref}
\newcommand{\sref}[2]{\hyperlink{source:#1:#2}{[S#2]}}
\setlength{\emergencystretch}{3em}
\begin{document}
% problemId:30007738
\section*{Policies and completed fertility}\label{30007738}
\subsection{Definitions and mathematical statement}
Consider women aged twenty-five at baseline in one high-income country introducing a specified childcare subsidy. Let \(a=1\) denote eligibility for the subsidy through age forty and \(a=0\) the preceding childcare regime. Define cumulative live births by age \(s\) as \(F_{is}(a)\), for \(s\in\{30,35,40,45\}\), and discounted net public resource cost as \(C_i(a)\). The targets are
\[\tau_s=E[F_{is}(1)-F_{is}(0)],\qquad\rho=\frac{E[C_i(1)-C_i(0)]}{\tau_{45}},\]
where \(\rho\) is interpreted only when \(\tau_{45}>0\). The expectation includes all baseline eligible women, independently of later partnership or employment choices. Identify whether the policy changes completed fertility rather than merely advancing births, using randomized access, discontinuous eligibility, or a justified rollout design with explicit counterfactual trends. Account for childcare capacity, wage responses, migration, and simultaneous family policies. Follow-up before age forty-five requires transparent completion projections or bounds. Report employment and child outcomes separately, since additional births alone do not measure welfare. The fertility reviews support the topic, but this particular completed-fertility policy contrast was not verified as an explicitly stated open problem in readable primary text.

\par\textbf{Formulation and scope.} Proposed scoped research specification. The original catalogue prompt was a synthesis, and no canonical conjecture or verified original formulation is claimed. The supporting source, when retained, establishes context or existing results rather than this exact open question.

\par\textbf{Open-question provenance.} An explicit published open-question statement was not verified. Sources below are supporting literature and must not be treated as a first listing.

\par\textbf{Historical priority.} No readable primary passage explicitly stating this precise open question was verified. A supporting paper, abstract, or topic citation is insufficient to establish origin or unresolved status.

\subsection{Short English statement}
Determine whether a defined childcare subsidy increases births completed by age forty-five rather than merely changing birth timing, and measure its resource cost.

\subsection{Sources}
\begin{itemize}\raggedright
\item[\textnormal{[S1]}]\hypertarget{source:30007738:1}{} Melissa S. Kearney, Phillip B. Levine. 2026. Why Is Fertility So Low in High-Income Countries?. Abstract.
\url{https://www.aeaweb.org/articles?id=10.1257/jel.20261786}.
DOI: 10.1257/jel.20261786.
{\small\textit{Supporting or current literature; see evidence qualification. The published review discusses explanations for falling fertility; the proposed counterfactual decomposition and completed-fertility policy contrast are our specifications.}}
\end{itemize}
\end{document}
